% !TeX root = ../thuthesis-example.tex

\nomenclature[00-delta-pes]{$\Delta$PES}{跨层级必需性评分差异}
\nomenclature[aa]{aa}{氨基酸 (Amino Acid)}
\nomenclature[attention]{Attention}{注意力机制 (Attention Mechanism)}
\nomenclature[atp6v1b2]{ATP6V1B2}{V-ATPase V1结构域B2亚基}
\nomenclature[auprc]{AUPRC}{精准率-召回率曲线下面积 (Area Under the Precision-Recall Curve)}
\nomenclature[auroc]{AUROC}{样本特征曲线下面积 (Area Under the Receiver Operating Characteristic Curve)}
\nomenclature[bert]{BERT}{基于Transformer的双向编码器表征 (Bidirectional Encoder Representations from Transformers)}
\nomenclature[casp14]{CASP14}{第14届蛋白质结构预测关键评测 (Critical Assessment of Structure Prediction, 14th edition)}
\nomenclature[ceres]{CERES}{拷贝数效应校正模型}
\nomenclature[cnn]{CNN}{卷积神经网络 (Convolutional Neural Network)}
\nomenclature[crispr-cas9]{CRISPR-Cas9}{成簇规律间隔短回文重复序列及其相关蛋白9 (Clustered Regularly Interspaced Short Palindromic Repeats-CRISPR associated protein 9)}
\nomenclature[depmap]{DepMap}{癌症依赖图谱项目 (Cancer Dependency Map)}
\nomenclature[dtag]{dTAG}{靶向降解标签系统 (degradation TAG system)}
\nomenclature[ecm]{ECM}{细胞外基质 (Extracellular Matrix)}
\nomenclature[esm]{ESM}{进化尺度建模蛋白语言模型 (Evolutionary Scale Modeling)}
\nomenclature[exac]{ExAC}{外显子组聚合联盟 (Exome Aggregation Consortium)}
\nomenclature[gnomad]{gnomAD}{人类基因组聚合数据库 (Genome Aggregation Database)}
\nomenclature[gnn]{GNN}{图神经网络 (Graph Neural Network)}
\nomenclature[go]{GO}{基因本体论 (Gene Ontology)}
\nomenclature[hl-60]{HL-60}{人早幼粒细胞白血病细胞系 (Human Promyelocytic Leukemia Cell Line)}
\nomenclature[loeuf]{LOEUF}{基因功能缺失观察值/期望值比例置信上界 (Loss-of-function Observed/Expected Upper bound Fraction)}
\nomenclature[lof]{LoF}{功能缺失变异 (Loss of Function)}
\nomenclature[lstm]{LSTM}{长短期记忆网络 (Long Short-Term Memory)}
\nomenclature[lylap]{LyLAP}{溶酶体亮氨酸氨基肽酶 (Lysosomal Leucine Aminopeptidase)}
\nomenclature[mageck]{MAGeCK}{基于模型的全基因组CRISPR-Cas9敲除分析方法 (Model-based Analysis of Genome-wide CRISPR-Cas9 Knockout)}
\nomenclature[mlm]{MLM}{掩码语言建模 (Masked Language Modeling)}
\nomenclature[mlp]{MLP}{多层感知机 (Multilayer Perceptron)}
\nomenclature[msa]{MSA}{多序列比对 (Multiple Sequence Alignment)}
\nomenclature[mw]{MW}{分子量 (Molecular Weight)}
\nomenclature[nadph]{NADPH}{还原型烟酰胺腺嘌呤二核苷酸磷酸 (Reduced Nicotinamide Adenine Dinucleotide Phosphate)}
\nomenclature[nets]{NETs}{中性粒细胞胞外诱捕网 (Neutrophil Extracellular Traps)}
\nomenclature[pes]{PES}{蛋白质必需性评分 (Protein Essentiality Score)}
\nomenclature[pi]{pI}{等电点 (Isoelectric Point)}
\nomenclature[pic]{PIC}{蛋白重要性计算器 (Protein Importance Calculator)}
\nomenclature[plbd1]{PLBD1}{磷脂酶B结构域包含蛋白1 (Phospholipase B Domain-containing 1)}
\nomenclature[plm]{PLM}{蛋白语言模型 (Protein Language Model)}
\nomenclature[pma]{PMA}{佛波醇-12-肉豆蔻酸-13-乙酸酯 (Phorbol 12-Myristate 13-Acetate)}
\nomenclature[ppi]{PPI}{蛋白质相互作用网络 (Protein-Protein Interaction)}
\nomenclature[prott5]{ProtT5}{基于T5架构的预训练蛋白语言模型}
\nomenclature[protac]{PROTAC}{蛋白水解靶向嵌合体 (Proteolysis Targeting Chimera)}
\nomenclature[pssm]{PSSM}{位置特异性评分矩阵 (Position-Specific Scoring Matrix)}
\nomenclature[rnn]{RNN}{循环神经网络 (Recurrent Neural Network)}
\nomenclature[ros]{ROS}{活性氧 (Reactive Oxygen Species)}
\nomenclature[sgrna]{sgRNA}{单导RNA (single-guide RNA)}
\nomenclature[shrna]{shRNA}{短发夹RNA (short hairpin RNA)}
\nomenclature[t5]{T5}{文本到文本迁移Transformer (Text-to-Text Transfer Transformer)}
\nomenclature[tape]{TAPE}{蛋白嵌入任务评估基准 (Tasks Assessing Protein Embeddings)}
\nomenclature[v-atpase]{V-ATPase}{液泡型ATP酶 (Vacuolar-type ATPase)}

\printnomenclature[2.5cm]


% 也可以使用nomencl宏包，需要在导言区
% \usepackage{nomencl}
% \makenomenclature

% 在这里输出符号说明
% \printnomenclature[3cm]

% 在正文中的任意为都可以标题
% \nomenclature{PI}{聚酰亚胺}
% \nomenclature{MPI}{聚酰亚胺模型化合物，N-苯基邻苯酰亚胺}
% \nomenclature{PBI}{聚苯并咪唑}
% \nomenclature{MPBI}{聚苯并咪唑模型化合物，N-苯基苯并咪唑}
% \nomenclature{PY}{聚吡咙}
% \nomenclature{PMDA-BDA}{均苯四酸二酐与联苯四胺合成的聚吡咙薄膜}
% \nomenclature{MPY}{聚吡咙模型化合物}
% \nomenclature{As-PPT}{聚苯基不对称三嗪}
% \nomenclature{MAsPPT}{聚苯基不对称三嗪单模型化合物，3,5,6-三苯基-1,2,4-三嗪}
% \nomenclature{DMAsPPT}{聚苯基不对称三嗪双模型化合物（水解实验模型化合物）}
% \nomenclature{S-PPT}{聚苯基对称三嗪}
% \nomenclature{MSPPT}{聚苯基对称三嗪模型化合物，2,4,6-三苯基-1,3,5-三嗪}
% \nomenclature{PPQ}{聚苯基喹噁啉}
% \nomenclature{MPPQ}{聚苯基喹噁啉模型化合物，3,4-二苯基苯并二嗪}
% \nomenclature{HMPI}{聚酰亚胺模型化合物的质子化产物}
% \nomenclature{HMPY}{聚吡咙模型化合物的质子化产物}
% \nomenclature{HMPBI}{聚苯并咪唑模型化合物的质子化产物}
% \nomenclature{HMAsPPT}{聚苯基不对称三嗪模型化合物的质子化产物}
% \nomenclature{HMSPPT}{聚苯基对称三嗪模型化合物的质子化产物}
% \nomenclature{HMPPQ}{聚苯基喹噁啉模型化合物的质子化产物}
% \nomenclature{PDT}{热分解温度}
% \nomenclature{HPLC}{高效液相色谱（High Performance Liquid Chromatography）}
% \nomenclature{HPCE}{高效毛细管电泳色谱（High Performance Capillary lectrophoresis）}
% \nomenclature{LC-MS}{液相色谱-质谱联用（Liquid chromatography-Mass Spectrum）}
% \nomenclature{TIC}{总离子浓度（Total Ion Content）}
% \nomenclature{\textit{ab initio}}{基于第一原理的量子化学计算方法，常称从头算法}
% \nomenclature{DFT}{密度泛函理论（Density Functional Theory）}
% \nomenclature{$E_a$}{化学反应的活化能（Activation Energy）}
% \nomenclature{ZPE}{零点振动能（Zero Vibration Energy）}
% \nomenclature{PES}{势能面（Potential Energy Surface）}
% \nomenclature{TS}{过渡态（Transition State）}
% \nomenclature{TST}{过渡态理论（Transition State Theory）}
% \nomenclature{$\increment G^\neq$}{活化自由能（Activation Free Energy）}
% \nomenclature{$\kappa$}{传输系数（Transmission Coefficient）}
% \nomenclature{IRC}{内禀反应坐标（Intrinsic Reaction Coordinates）}
% \nomenclature{$\nu_i$}{虚频（Imaginary Frequency）}
% \nomenclature{ONIOM}{分层算法（Our own N-layered Integrated molecular Orbital and molecular Mechanics）}
% \nomenclature{SCF}{自洽场（Self-Consistent Field）}
% \nomenclature{SCRF}{自洽反应场（Self-Consistent Reaction Field）}
